%%%%%%%%%%  scan idx 78  |  folio 65  %%%%%%%%%%
tout cas facilement la deuxième formule (2.35) modulo le sous-groupe de torsion de $G(\underline K(X))$, et en caractéristique $0$, cette même formule est une conséquence du théorème de Riemann--Roch. L'intérêt de (2.35) (lorsqu'elle est vraie) tient à ce qu'elle montre que $\varphi$ est injective modulo des groupes de torsion (peut-être toujours nuls~!\footnote{Non, cf.\ note de la page 30.}), de sorte qu'on «~ne perd pas beaucoup~» à définir les classes de Chern comme des éléments de $G(\underline K(X))$~; si l'on fait le produit tensoriel par $\Q$, comme on le fera plus loin pour le théorème de Riemann--Roch, les deux définitions des classes de Chern seraient alors identiques.

Soit toujours $X$ quasi-projective non singulière~; soit de plus $Y$ une sous-variété non singulière de $X$ et soit $i$ l'injection de $Y$ dans $X$. On a alors, pour $i^!$ et $i_!$ les formules suivantes~:
\begin{enumerate}[label=(\roman*)]
\item $i^!$ est un homomorphisme de $\lambda$-anneaux
\item $i_!(x.i^!(y))=i_!(x).y$, d'où en particulier $i_!i^!(y)=\xi y$ avec $\xi=i_!i^!(1)=\gamma_X(Y)$,
\item $i^!(i_!(x))=x.\lambda_{-1}(N)$, d'où en particulier $i^!(\xi)=i^!i_!(1)=\lambda_{-1}(N)$,
\item $i_!(x)i_!(x')=i_!(xx'\lambda_{-1}(N))$,
\end{enumerate}
où, par définition, $N$ est la classe dans $\underline K(Y)$ du fibré conormal à $Y$ dans $X$. Les formules (i) et (ii) sont vraies pour des morphismes quelconques, non nécessairement des injections, et (iii) et (iv) se vérifient sur la définition (2.15) et la formule analogue pour les images inverses (en prenant pour $x$ et $x'$ des classes de fibrés vectoriels, et en se servant du fait que
\begin{equation}\tag{2.36}
 \Tor^{\cO_X}_i(\cO_Y,\cO_Y)=\cO_Y(\wedge^i N).
\end{equation}
)
